\abstract{РЕФЕРАТ}

Объем работы равен \formbytotal{lastpage}{страниц}{е}{ам}{ам}. Работа содержит \formbytotal{figurecnt}{иллюстраци}{ю}{и}{й}, \formbytotal{tablecnt}{таблиц}{у}{ы}{}, \arabic{bibcount} библиографических источников и \formbytotal{числоПлакатов}{лист}{}{а}{ов} графического материала. Количество приложений – 2. Графический материал представлен в приложении А. Фрагменты исходного кода представлены в приложении Б.

Перечень ключевых слов: интерфейс, клиентская сторона, элементы, кнопки, меню, поля ввода, иконки, визуальная часть, стиль, графический интерфейс, текст, вайрфрейм, функциональность, триггер, уведомления, мокап, логотип, компоненты, виджеты, фирменный стиль, шапка, хеддер, тело сайта, боди, чекбокс, диалог, сообщения, логин, пароль, пользователь, логика взаимодействия, разметка, роли, доступность, переходы, анимации, тестирование, адаптивный дизайн, медиазапросы, брейкпоинты, сессия, авторизация, регистрация, аутентификация.

Объектом разработки является клиентская сторона закрытого корпоративного мессенджера. Элементы, которые видит пользователь и с которыми взаимодействует при использовании веб-приложения. Объект разработки ввключает в себя визуальную состовляющую (дизайн) и функциональную (механизмы взаимодействия), направленные на создание удобного, интуитивно понятного и визуально привлекательного опытя для пользователей, который позволяет им легко общаться и взаимодействовать с другими людьми. 

Цель выпускной квалификационной работы заключается в создании эффективной, удобной, простой и интуитивно понятной среды для внутренней коммуникации и совместной работы сотрудников, которая обеспечит и поспособствует повышению продуктивности и достижению бизнес-целей компании.

При разработке системы использовалась собственная архитектура с использованием клиент-серверной модели взаимодействия. Клиентская сторона, реализованная на HTML, CSS и JavaScript, отвечает за формирование пользовательского интерфейса и обработку взаимодействия на стороне пользователя.

\selectlanguage{english}
\abstract{ABSTRACT}
  
The volume of work is \formbytotal{lastpage}{page}{}{s}{s}. The work contains \formbytotal{figurecnt}{illustration}{}{s}{s}, \formbytotal{tablecnt}{table}{}{s}{s}, \arabic{bibcount} bibliographic sources and \formbytotal{числоПлакатов}{sheet}{}{s}{s} of graphic material. The number of applications is 2. The graphic material is presented in annex A. The layout of the site, including the connection of components, is presented in annex B.

List of Keywords: interface, client-side, elements, buttons, menu, input fields, icons, visual part, style, graphical user interface, text, wireframe, functionality, trigger, notifications, mockup, logo, components, widgets, corporate identity, header, body (of website), checkbox, dialog, messages, login, password, user, interaction logic, layout, roles, accessibility, transitions, animations, testing, adaptive design, media queries, breakpoints, session, authorization, registration, authentication.

The object of development is the client-side of a closed corporate messenger. These are the elements that the user sees and interacts with when using the web application. The object of development includes both the visual component (design) and the functional component (interaction mechanisms), aimed at creating a convenient, intuitive, and visually appealing experience for users, allowing them to easily communicate and interact with other people.

The goal of the final qualifying work is to create an effective, convenient, simple, and intuitive environment for internal communication and employee collaboration, which will ensure and contribute to increasing productivity and achieving the company's business goals.

The system was developed using a custom architecture based on a client-server interaction model. The client-side, implemented using HTML, CSS, and JavaScript, is responsible for forming the user interface and processing user-side interactions.
\selectlanguage{russian}
